\documentclass[11pt,letterpaper]{article}

\usepackage[margin=1in]{geometry}
\usepackage{amsmath,amssymb,amsthm}
\usepackage{booktabs}
\usepackage{graphicx}
\usepackage{xcolor}
\usepackage[hidelinks]{hyperref}
\usepackage{microtype}
\usepackage{enumitem}
\usepackage{titlesec}
\usepackage{caption}

\definecolor{mDarkTeal}{HTML}{1c3d6b}
\definecolor{mLightBrown}{HTML}{2f6fb3}

\titleformat{\section}{\large\bfseries\color{mDarkTeal}}{\thesection}{0.6em}{}
\titleformat{\subsection}{\normalsize\bfseries}{\thesubsection}{0.5em}{}
\titlespacing*{\section}{0pt}{1.4em}{0.6em}
\titlespacing*{\subsection}{0pt}{1.0em}{0.4em}

\hypersetup{colorlinks=true, linkcolor=mDarkTeal, citecolor=mDarkTeal, urlcolor=mLightBrown}

\newcommand{\R}{\mathbb{R}}
\newcommand{\Gop}{\mathcal{G}}
\newcommand{\sgrad}{\operatorname{sg}}
\newcommand{\relL}[1]{\mathrm{rel\text{-}L^{#1}}}

\title{\textbf{Improving Long-Rollout Accuracy of the DNO Surrogate}\\[2pt]
       \large A status report on training, diagnostics, and architectural directions}
\author{Jonathan Ma}
\date{4 June 2026}

\begin{document}
\maketitle

\begin{abstract}
\noindent
We report on a focused effort to close the long-rollout accuracy gap of a
1D Zakharov water-wave neural surrogate. Two training-time stability
methods have been implemented and partly evaluated: a one-step
pushforward loss (40-epoch run complete, an 80-epoch matched-compute
resume in progress) and a Hamiltonian-conservation loss (Modal training
in progress). Forensic analysis on the pushforward checkpoint vs.\ the
FNO baseline identifies a coherent, linear-in-$t$ in-band failure mode,
strongly predicted within each regime by initial-condition bandwidth and
maximum steepness ($\rho$ up to $+0.85$). A 1000-snapshot mmap'd audit
of the 73.7\,GB training dataset finds no data-coverage problem: every
test regime's initial conditions fall inside the training central 90\%
at the median. We summarise the empirical status of both training
methods, characterise the failure mode in terms of physical descriptors
of the initial condition, and rank concrete architectural and
training-loss changes by expected gain. The single largest unexploited
lever appears to be encoding the polynomial-in-$\eta$ structure of the
Craig--Sulem expansion into the architecture itself; longer-horizon
unrolled training with gradients through all steps is the most
straightforward training-side win.
\end{abstract}

% ===================================================================
\section{Background and scope}

The 1D Zakharov system on a periodic domain of length $L$ couples the
free-surface elevation $\eta(x,t)$ to the surface velocity potential
$\xi(x,t)$ via the Dirichlet--Neumann operator $\Gop(\eta;h)$:
\begin{align}
\partial_t \eta &= \Gop(\eta; h)\,\xi, \\
\partial_t \xi &= -g\,\eta - \tfrac{1}{2}\xi_x^{\,2}
                 + \tfrac{1}{2}\frac{\big(\Gop(\eta)\xi + \eta_x\,\xi_x\big)^2}{1 + \eta_x^{\,2}}.
\end{align}
Two neural surrogates are trained to approximate
$(\eta,\xi,h) \mapsto \Gop(\eta;h)\xi$ on single snapshots and then
iterated by an explicit time integrator at rollout time:
\begin{itemize}[leftmargin=1.3em,itemsep=2pt]
  \item \textbf{FNO baseline}, a 6-block FNO1d with width 128, modes 64,
        FiLM depth conditioning, soft-gating channel MLPs, and a
        linear-DNO baseline added in the output ($\sim 13.1$M params).
  \item \textbf{SpectralDNO}, a 6-block model with width 128, modes 64,
        latent 64, concatenating $\log h$ as a feature channel and
        producing the residual as a rank-$\mathrm{latent}$ spectral
        sandwich $C(\eta,h)^{\!\top}\!F(\eta,h)\,C(\eta,h)\,\xi$ added
        to the linear baseline ($\sim 17.0$M params). The output is
        \emph{linear in $\xi$ by construction}.
\end{itemize}

On the standard 10-regime evaluation suite (Tanaka solitons, Benjamin--
Feir breathers, Stokes waves, random sea, linear) both surrogates are
pointwise accurate but exhibit significant error growth over long
autoregressive rollouts. This report covers the work undertaken to
diagnose and mitigate the rollout error, and the recommended next steps.

% ===================================================================
\section{Training-time stability methods}

Both methods augment a baseline single-step Sobolev-$H^k$ loss with an
additional term targeting rollout stability. They are composable in the
training script but the currently in-flight runs use one or the other.

\subsection{Pushforward training}
\label{ssec:pf}

Given a data point $(\eta,\xi)$, advance the surrogate's own prediction
through $k$ Euler steps of the Zakharov system, then re-predict
$\Gop$ at the rolled-out state and penalise it against the analytical
Craig--Sulem expansion at order four:
\[
  \mathcal{L} \;=\; \mathcal{L}_{\text{single-step}}
  \;+\; \lambda_{\mathrm{pf}}\,\mathcal{L}\!\bigl(
        \Gop_\theta\!\bigl(\sgrad(\eta_k),\sgrad(\xi_k)\bigr),\;
        \Gop_{\text{CS}}(\eta_k,\xi_k)
        \bigr).
\]
Gradient flow is restricted to the final $\Gop$-prediction at the
rolled-out state; the trajectory and the analytical target are both
\texttt{stop\_gradient}-wrapped. The currently active configuration uses
$k=1$, $\lambda_{\mathrm{pf}}=1$, $\mathrm{d}t=0.01$, and matches the
model's depth saturation $h_{\max}=5$ when constructing the analytical
target.

\textbf{Intent.} Breaks the covariate-shift compounding of pure
single-step training: the surrogate must remain accurate on states that
it itself produces under iteration, not only on data-distribution states.

\textbf{Cost.} $(1+k)$ forward passes plus one vmapped Craig--Sulem
series evaluation per training step.

\subsection{Hamiltonian-conservation loss}
\label{ssec:hcons}

The Zakharov system is Hamiltonian with
\[
  H(\eta,\xi) \;=\; \tfrac{1}{2}\!\int_\Omega
    \bigl(\xi\,\Gop(\eta)\xi + g\,\eta^{2}\bigr)\,dx.
\]
An exact solution conserves $H$. The added loss penalises clipped
relative drift of $H$ over one Euler step of the surrogate's own
dynamics:
\[
  \mathcal{L}_H \;=\; \lambda_H\,
    w(s)\,\mathbb{E}\!\left[\,\min\!\left(
      \frac{|H_1 - H_0|}{|H_0|+\varepsilon},\;C\right)\right],
\]
with $H_0$ evaluated at the data state, $H_1$ at one Euler step, both
$\Gop$ values from the surrogate, $\sgrad$ on the advanced state, and a
linear warmup $w(s)=\min(s/s_{\mathrm{warm}},1)$. Configuration:
$\lambda_H=1$, $C=10$, $s_{\mathrm{warm}}=10\,000$,
$\mathrm{d}t=0.01$.

\textbf{Intent.} A physically motivated soft constraint: predictions
that violate energy conservation are penalised in proportion to how much
they violate it. Single-step horizon, so it does not stabilise long
rollouts directly the way pushforward does.

\textbf{Asymmetry.} The pushforward branch mirrors the model's
depth saturation when building the Craig--Sulem target; the H-cons
branch does not apply the clip when computing $H$ or the Zakharov RHS.
The relative-drift normalisation also makes the spatial measure
$\mathrm{d}x_{\mathrm{phys}}$ cancel, so the spatial integration factor
is effectively dead code in the current functional.

\subsection{Composition with model saturation}

The model itself clips $\log h$ at $\log 5$ so that $\tanh(hk)\to 1$ for
$k\geq 1$. The pushforward target respects this clip; the H-cons branch
does not. Either method can be disabled by setting its weight to zero.

% ===================================================================
\section{Evaluation framework}

\subsection{Regimes}
The eval suite covers ten regimes spanning depth (shallow/finite/deep),
spectral bandwidth (monochromatic/broadband), and physical character
(soliton/modulational/quasi-linear). Each provides 16 initial conditions
and a rollout horizon $T_{\max}$ matching the integrator's stable range.

\subsection{Linear-regime fix}
\label{ssec:linear-fix}

While diagnosing apparent surrogate failure on the ``linear'' regime, we
discovered that the regime's initial conditions are physically nonlinear
shallow-water cases ($\eta/h$ up to $0.30$, $kh\approx 0.6$) inherited
from the MATLAB legacy test set. The Zakharov solver used to produce
the ``ground truth'' was diverging on some of these. We replaced the
truth integrator with analytic per-mode Fourier propagation of the
\emph{linear} Zakharov system,
\[
\omega^2(k) = g\,k\,\tanh(hk),\qquad
\hat\eta(k,t) = \hat\eta_0(k)\cos\omega t
              + \frac{\Gop_{\mathrm{lin}}(k)}{\omega}\hat\xi_0(k)\sin\omega t,
\]
and analogously for $\hat\xi(t)$. Amplitude and top-$k$ structure are
preserved to better than $1\%$ over the full horizon. All linear-regime
comparisons in this report use the analytic truth.

\subsection{Metrics}
We use per-IC relative $L^{2}$ error of $\eta$,
$
\|\eta_{\mathrm{pred}}(t)-\eta_{\mathrm{truth}}(t)\|_2 /
\|\eta_{\mathrm{truth}}(t)\|_2,
$
reported at three time slices ($T_{\max}/10$, $T_{\max}/2$, $T_{\max}$)
or as full curves. ICs whose truth or prediction is non-finite anywhere
in the rollout are flagged and counted separately rather than imputed.

% ===================================================================
\section{Empirical findings}

\subsection{Pushforward training: a real but uneven win}

Comparing FNO baseline (80 epochs) to SpectralDNO + pushforward
(40 epochs), median rel-$L^{2}$ at the final time horizon:

\smallskip
\begin{center}
\small
\begin{tabular}{lrrrr}
\toprule
Regime & FNO med & DNO+PF med & $\Delta$ & NaN ICs (FNO / DNO+PF) \\
\midrule
tanaka\_g0          & 0.246  & 0.099  & $-60\%$       & 0 / 1 \\
tanaka\_g1          & 0.605  & 0.255  & $-58\%$       & 0 / 1 \\
bf\_g0              & 0.174  & 0.130  & $-25\%$       & 0 / 3 \\
bf\_g1              & 0.193  & 0.130  & $-33\%$       & 1 / 2 \\
bf\_modal           & 0.245  & 0.127  & $-48\%$       & 3 / 2 \\
linear (analytic)   & 1.170  & 0.341  & $-71\%$       & 0 / 4 \\
stokes\_deep        & 0.0036 & 0.0010 & $-72\%$       & 0 / 0 \\
stokes\_finite      & 0.0035 & 0.0056 & $+60\%$       & 0 / 0 \\
random\_sea\_deep   & 0.0041 & 0.120  & $+2799\%$     & 0 / 0 \\
random\_sea\_finite & 0.0042 & 0.091  & $+2088\%$     & 0 / 0 \\
\bottomrule
\end{tabular}
\end{center}
\smallskip

Pushforward training delivers substantial slope reductions (5--25$\times$)
on long-horizon, broadband, shallow-water regimes
(Tanaka, BF, modal, linear), but \emph{regresses substantially} on
random-sea regimes where the FNO baseline was already near-zero. NaN
rates also rise under DNO+PF in some shallow regimes (e.g.\ 4 of 16
linear ICs blow up under DNO+PF; 0 under FNO).

\textbf{Confound.} The DNO+PF run is 40 epochs vs.\ 80 for the FNO
baseline, with an architecture that has $\sim 30\%$ more parameters.
A matched-compute 80-epoch DNO+PF resume run (with reset Adam moments
because of an optax pytree-shape mismatch on the saved opt-state) is in
progress as of the cutoff for this report and will form the
apples-to-apples comparison.

\subsection{Failure mode is coherent in-band, not high-$k$}
\label{ssec:failure-mode}

Across the three worst finite-error ICs (one each from linear, bf\_modal,
bf\_g1) the error grows linearly or with slight super-linearity in
time, $r^{2}_{\mathrm{lin}}$ between $0.53$ and $0.93$. Decomposing the
error spectrum into a peak band (centered on the truth's dominant
wavenumber) versus a high-$k$ band, the per-IC cosine similarity
between the error spectrum and the truth spectrum is $0.57$--$1.00$ for
FNO, with high-$k$ Gibbs-style energy fraction below 0.18 in every
worst case. NaN blow-ups in bf\_g1 and bf\_modal occur at $t = 81$--$189$
out of $T_{\max}=200$ while $|\eta|_{\max}$ stays small ($0.04$--$0.12$),
i.e.\ numerical divergence after extended coherent drift, not amplitude
blowup.

\textbf{Specifically.} On the linear regime, FNO at $t=2$ has $99\%$ of
its error power in sideband modes $k=10$--$50$ while the truth has all
its power at $k=6$. This is sideband hallucination, and is the
architectural cost of FNO mixing $\xi$ into nonlinear features. By
construction the SpectralDNO cannot generate this kind of error on the
linear regime --- consistent with its substantial improvement there.

\subsection{Physics descriptors that predict failure}

Within-regime Spearman correlations between IC physical descriptors and
final-time rel-$L^{2}$:
\begin{itemize}[leftmargin=1.3em,itemsep=2pt]
  \item Tanaka g0/g1: spectral \textbf{bandwidth} $\rho = +0.76\,/\,+0.78$
        for FNO, $+0.84\,/\,+0.85$ for DNO+PF --- strongest signal in the
        dataset.
  \item Random-sea (deep): \textbf{max steepness} $\rho = +0.63$ (FNO)
        $/\,+0.84$ (DNO+PF).
  \item Stokes (finite): \textbf{max steepness} $\rho = +0.70\,/\,+0.83$.
  \item BF g0/g1: \textbf{$k_{\mathrm{peak}}$} $\rho \in [+0.33,+0.87]$.
\end{itemize}
Global correlations with $kh$ and depth ($\rho \approx -0.4$) appear
strong but are confounded --- Tanaka has both shallow water and broad
spectra, so the marginal global view conflates two distinct effects.

\textbf{Read-off.} The single most informative descriptor of where the
surrogate fails is \emph{not} a distributional one (kh, depth) but a
\emph{structural} one (bandwidth, steepness). High-steepness and
broadband ICs need the nonlinear-in-$\eta$ part of $\Gop$ to be
accurate, which is exactly the part of the architecture we suspect is
under-parameterised.

\subsection{Rollout error vs.\ initial wavenumber}

The headline diagnostic plot (Fig.~\ref{fig:err-vs-k}) takes three ICs
per regime binned by initial mean wavenumber and plots, on shared time
axes, the relative $L^{2}$ error growth and the mean wavenumber of the
truth, $\langle k\rangle(t) = \sum_k k\,|\hat\eta(k,t)|^{2}/\sum_k
|\hat\eta(k,t)|^{2}$.

\begin{figure}[t]
\centering
\includegraphics[width=0.9\linewidth]{../../outputs/error_vs_k/tanaka_g0_err_vs_k.png}
\caption{Tanaka g0: error growth and truth mean wavenumber vs.\ time
for three ICs binned by initial mean wavenumber (red = low-$k$,
blue = mid-$k$, green = high-$k$). Top: log-scale rel-$L^{2}$ $\eta$-error.
Bottom: $\langle k\rangle(t)$ of the truth. The wave's mean wavenumber
is essentially stationary; the rollout error grows by roughly two orders
of magnitude over the horizon, more rapidly for higher-$k$ ICs.}
\label{fig:err-vs-k}
\end{figure}

The pattern visible in Tanaka g0 holds qualitatively in most regimes:
higher-$k$ ICs produce larger and faster-growing rollout error, while
the truth's mean wavenumber is essentially stationary (for the linear
regime, exactly so). This rules out spectral drift of the underlying
wave as an explanation for the error growth.

% ===================================================================
\section{Data quality audit}

A 1000-snapshot mmap'd sample of the 73.7\,GB
\texttt{combined\_dataset\_v3.npz} found no NaN/inf in over $10^{6}$
sampled voxels of $\eta$, $\xi$, or $\Gop\xi$. Every test regime's IC
distribution falls inside the training central 90\% at the median
for amplitude, steepness, and dominant wavenumber. Two tails are worth
noting:
\begin{itemize}[leftmargin=1.3em,itemsep=2pt]
  \item The maximum-amplitude Tanaka cases nip the training $p99$ in
        $\eta$ and $\xi$ amplitude.
  \item The \texttt{stokes\_finite} dominant wavenumber sits at the
        upper edge of training $k$ ($\sim 26$); modes are still within
        the network's $\mathrm{modes}=64$ cutoff.
\end{itemize}
One eval-side issue surfaced: one of the 16 \texttt{bf\_g1} ICs has the
\emph{truth} integrator diverge to NaN at $t\approx 146$ of $T_{\max}=200$.
This pollutes any aggregate metric on that regime. It is not a
training-data problem.

\textbf{Conclusion.} The rollout-accuracy ceiling is not a
data-coverage problem. Improvements have to come from the
model or the training-time loss, not from acquiring more data.

% ===================================================================
\section{Architectural analysis}

Both surrogates output
\[
  \Gop_\theta(\eta;h)\xi = G_{0}(h)\,\xi + R_{\theta}(\eta,h)\,\xi,
\]
where $G_{0}(h) = |D|\tanh(h|D|)$ is the flat-surface linear DNO
computed in closed form and $R_{\theta}$ is the learned residual.
\begin{itemize}[leftmargin=1.3em,itemsep=2pt]
  \item In FNO, $R_{\theta}$ is the standard FNO trunk applied to
        $(\eta,\xi,\partial\eta,\partial\xi,\dots)$ with FiLM depth
        conditioning; the trunk mixes $\xi$ into the nonlinear
        features.
  \item In SpectralDNO, $R_{\theta} = C^{\!\top}(\eta,h)\,F(\eta,h)\,
        C(\eta,h)$ is a rank-$\mathrm{latent}$ spectral sandwich. By
        construction $R_\theta$ is \emph{linear in $\xi$}, and the
        spectral filter $F(k)$ is constructed from a \emph{mean-pooled}
        trunk feature --- so the filter is one shape per IC, not
        spatially varying.
\end{itemize}

Two architectural observations follow directly from \S\ref{ssec:failure-mode}:

\textbf{Linearity in $\xi$ is a real win on the linear regime.} FNO's
sideband hallucination on the linear regime is a structural
consequence of mixing $\xi$ into nonlinear features. SpectralDNO cannot
produce this kind of error, and indeed dominates FNO on the linear
regime by $-71\%$ in median rel-$L^{2}$.

\textbf{The polynomial-in-$\eta$ structure is under-parameterised.}
The Craig--Sulem expansion
$\Gop(\eta)\xi = G_{0}(h)\xi + G_{1}(\eta,h)\xi + G_{2}(\eta\otimes
\eta,h)\xi + \cdots$
has $G_{n}$ multilinear in $\eta^{\otimes n}$ and linear in $\xi$.
SpectralDNO's single learned sandwich is shape-compatible with one
$G_{1}$ contribution; the architecture has no mechanism to express
genuine $G_{2}$ or higher contributions, which are exactly the terms
that dominate when steepness or bandwidth is large.

% ===================================================================
\section{Recommended directions, ranked}

We rank by expected gain $\times$ implementation effort, and
distinguish architectural changes from training-loss changes.

\subsection{Architecture}

\paragraph{1.\;Polynomial-in-$\eta$ block structure (Craig--Sulem-shaped).}
Replace the single learned sandwich with a stack of blocks each shaped
like a learned $G_{n}$:
\[
  \text{DNOBlock}_{n}(\eta,\xi,h)
  \;\sim\; M_{\mathrm{out}}(D,h)\,
  \sum_{i}\bigl[\Phi_{i}(\eta)\cdot M_{i}(D,h)\,\xi\bigr],
\]
where $\Phi_{i}$ are learned spatial features of $\eta$ (e.g.\ pointwise
polynomials or small spatial MLPs) and $M_{i}(D,h)$ are learned
depth-conditioned Fourier multipliers. Sum the outputs of 3--5 blocks
and add to $G_{0}\xi$. The whole architecture is then \emph{linear in
$\xi$ by construction} (preserving the win on the linear regime) and
\emph{polynomial in $\eta$ by construction} (addressing the
high-steepness failure mode). Theoretical grounding: this \emph{is} the
Craig--Sulem series with learned multipliers. Estimated implementation:
$\sim$150 lines of model code; no training-loop changes.

\paragraph{2.\;Depth-and-$k$-aware spectral filter.}
The current spectral filter is
$F = \mathrm{softplus}(\mathrm{MLP}(\mathrm{mean\_pool}(\text{trunk}))$,
with no explicit access to $k$ or $h$. Replace with
\[
  F(k,h,\text{features}) = \mathrm{softplus}\!\bigl(\mathrm{MLP}(k,\,h,\,
    \tanh(hk),\, k\tanh(hk),\,\text{features})\bigr).
\]
The first four inputs are the natural features of the dispersion
relation; the MLP only needs to learn perturbations on this ansatz.
Probably the highest gain-per-line change available; $\sim$20 lines.

\paragraph{3.\;Spectral derivatives and half-derivative features.}
Replace the finite-difference $\partial_{x},\partial_{xx}\eta$ features
with spectral derivatives, and add $|D|^{1/2}\eta$ and the Hilbert
transform $\mathcal{H}\eta$ as input features. $|D|^{1/2}$ is the
natural scaling that appears in $G_{1}$ for deep water; currently the
network must reconstruct it from local stencils. $\sim$30 lines, no
parameter increase.

\paragraph{Smaller wins.}
Smooth depth saturation ($h_{\mathrm{eff}} = h_{\max}\tanh(h/h_{\max})$
rather than hard clip); port FNO's per-mode weight-norm and learnable
gain to SpectralDNO's spectral conv; allow the spectral filter to be
spatially varying via attention pooling rather than mean pooling.

\subsection{Training-time}

\paragraph{4.\;K-step unrolled training with gradients through all steps,
$k=4$--$8$, with curriculum.}
List, Chen, Bali \& Thuerey~(\href{https://arxiv.org/abs/2402.12971}{arXiv:2402.12971})
show that stop-gradient pushforward has a bias scaling like $k^{2}$
relative to the true multi-step gradient, and that $k\leq 4$ is often
insufficient. Replace
\(
\text{loss} = \mathrm{MSE}(\text{model}(\sgrad(\text{model}(u))), \text{target})
\)
with
\(
\text{for } i\in[1,k]:\; u \leftarrow \text{model}(u);\;
\text{loss}\,{+}{=}\,w_{i}\,\mathrm{MSE}(u, \text{target}_{i})
\),
no stop-gradient, and a curriculum that ramps $k$ from 1 to 8 over
training. Memory scales linearly in $k$; at width 128, modes 64, $k=4$
is comfortable on a single H100 with halved batch size.

\paragraph{5.\;Predict tendency, integrate externally with RK4.}
Bryutkin \& Karniadakis~(\href{https://arxiv.org/abs/2412.13074}{arXiv:2412.13074})
report large rollout improvements from reformulating the target as
$\partial_{t}u$ and using a real RK4 step at inference. Strictly
speaking the DNO output is already a tendency (it is
$\partial_{t}\eta$), so the analogous reformulation would be on the
$\xi$ side: have the model also predict the Bernoulli RHS rather than
relying on a hand-coded formula at rollout. This decouples the
surrogate from carrying the dispersive phase rotation across one large
step. Modest implementation effort; high theoretical motivation.

\paragraph{6.\;PDE-Refiner-style noise-conditioned refinement.}
Lippe et al.~(\href{https://arxiv.org/abs/2308.05732}{arXiv:2308.05732}):
standard MSE training neglects low-amplitude high-$k$ modes, exactly
the modes that drive modulational instability in Zakharov. Replace
one-shot next-step prediction with $K=3$ iterative noise-conditioned
denoising passes. Inference cost $K\times$; reported gain is smaller
on FNO than U-Net so the SpectralDNO is a better candidate here.

\subsection{Skip list}

Strict-symplectic / hard-Hamiltonian architectures (wrong constraint
level for this discretisation), global Lipschitz bounds on the operator
(kills the modulational dynamics for dispersive PDEs), simple capacity
scaling (current architecture is 17M params vs.\ FNO's 13M; capacity is
not the bottleneck). See the literature review
(\texttt{notes/rollout\_accuracy/literature\_review.md}) for full
reasoning on skipped directions.

% ===================================================================
\section{Open questions}

\begin{itemize}[leftmargin=1.3em,itemsep=3pt]
  \item Does the random-sea regression under DNO+PF survive the 80-epoch
        matched-compute resume, or is it a training-budget artifact?
  \item Is the H-conservation loss as currently written --- evaluating
        $H_{0}$ with the \emph{surrogate's} $\Gop$-prediction rather
        than the data label --- the intended functional, or should
        $H_{0}$ use the true $\Gop\xi$ label?
  \item Does a Craig--Sulem-shaped architecture without pushforward
        already beat the current pushforward DNO? If so, the polynomial
        structure is the operative inductive bias, not the training
        loss.
  \item Per-IC NaN traces in bf\_g1/modal: does divergence originate in
        the $\Gop$-prediction or in the $\xi$-side Bernoulli RHS?
        Knowing this changes which loss term to weight.
\end{itemize}

% ===================================================================
\section{Status of in-flight work}

\begin{itemize}[leftmargin=1.3em,itemsep=2pt]
  \item \textbf{DNO+PF 80-epoch resume:} relaunched after an optax
        opt-state pytree-shape mismatch on the first attempt; running
        with \texttt{--reset\_opt\_state} (Adam moments rewarm in
        $\sim$100 steps, negligible against 40 epochs). As of writing,
        epoch 2/40 complete, val\_loss 0.00566. Rewarm is consistent
        with the original best 0.00490 at epoch 31.
  \item \textbf{H-cons training (Modal, L40S):} launched with linear
        warmup and per-sample clip; first evaluation pending.
  \item \textbf{Eval suite:} run on FNO baseline and DNO+PF
        40-epoch; all 10 regimes evaluated, trajectories stored under
        \texttt{outputs/<run>/eval\_suite/*\_trajs.npz}. Will be re-run
        on the resumed and H-cons checkpoints when they finish.
\end{itemize}

% ===================================================================
\section*{Artifacts}

\begin{itemize}[leftmargin=1.3em,itemsep=1pt]
  \item Data audit: \texttt{notes/rollout\_accuracy/data\_audit\_report.md}
        and \texttt{data\_audit.json}.
  \item Rollout forensics:
        \texttt{notes/rollout\_accuracy/rollout\_forensics\_report.md}
        and \texttt{rollout\_forensics.json}. The exploratory analysis scripts
        remain available through Git history.
  \item Literature scan:
        \texttt{notes/rollout\_accuracy/literature\_review.md}.
  \item Per-regime error-vs-wavenumber plots:
        \texttt{outputs/error\_vs\_k/*\_err\_vs\_k.png}; the historical site
        builder remains available through Git history, and the deployed site is
        at \href{https://stunning-cocada-db3a7d.netlify.app}{stunning-cocada-db3a7d.netlify.app}.
  \item Linear-regime fix:
        \texttt{solver/evals/eval\_suite.py}, function
        \texttt{truth\_rollout\_linear\_analytic}.
  \item Pushforward and H-cons training implementation:
        \texttt{train-jax-10m/1d\_dno\_fno\_jax.py}, lines 396--463
        (loss assembly), 344--367 (Zakharov Euler step and analytical
        target wrapper).
  \item Companion slides:
        \texttt{notes/rollout\_stability\_approaches.tex/.pdf}.
\end{itemize}

\end{document}
